\section{产品定义：利率暴露，不是保护性现金}

511010（国泰上证5年期国债ETF）由国泰基金管理有限公司管理、中国建设银行托管，以优化抽样跟踪上证5年期国债指数收益率。产品资料概要披露，建仓完成后不低于 90\% 的基金资产应投资于标的指数成份券或备选成份券；管理费 0.15\%，托管费 0.05\%，年化运营费用估算 0.21\%。因此其经济本质是\textbf{集中于五年期国债曲线的 ETF 利率敞口}，不是到期本金确定的存款，也不是没有净值波动的“保护性流动性”。

\begin{exhibitbox}[表 3：511010 的已核验事实与未完成模块]
\small
\begin{tabularx}{\exhibitboxwidth}{L{3.4cm}X}
\toprule
\textbf{模块} & \textbf{已核验事实 / 结论} \\
\midrule
2025 年报（经审计） & 基金资产净值 CNY3.021bn；每份 NAV CNY140.027；年度 NAV 回报 0.59\%；现金分配 CNY21.042/10 份，合计 CNY38.003m。历史分配不等于未来收益。\\
2026Q2（未经审计） & 2026-06-30 NAV CNY5.117bn、每份 NAV CNY140.753、三个月 NAV 回报 0.78\%；国债占 NAV 98.89\%。\\
穿透集中度 & Q2 前五项国债占 NAV 87.30\%，最大单项“25附息国债20”占 20.75\%。这是成份券集中度描述，不是久期风险模型。\\
交易所观察 & 2026-08-12 16:29:05 最后价 CNY141.084，当日成交额 CNY1.345bn。快照没有同日 IOPV/NAV、点差、深度或容量。\\
估值状态 & 价格相对 2026-06-30 NAV 的日期错配诊断约 +0.24\%；相隔 43 天，不能称为当前溢价或公允价值。\\
行动 & \qzero；任何条件性 15\% 参考权重都未获批准。\\
\bottomrule
\end{tabularx}
\sourcenote{年报、Q2、产品资料概要、持仓表和上交所快照：\sourceid{S04--S06, S27, S28}。}
\end{exhibitbox}

\section{真正的风险变量}

该产品的风险问题不是“国债是否安全”这一抽象判断，而是：当收益率曲线平移或陡峭化时，基金的修正久期、DV01 和凸性是多少；优化抽样相对指数如何偏离；当前成交价格是否可以以给定规模执行；以及账户、税费与申赎机制是否合规。现有源包没有这些当前字段。即使披露资产几乎全为国债，也不能把它改写为无风险现金或替代未命名的保护性流动性类别。

\begin{riskbox}[升级触发条件]
在 q 能从零变动之前，需要同一时点的基金 NAV/IOPV、完整成份券与修正久期/DV01/凸性、至少 20 个交易日的点差/深度/冲击成本、申赎状态以及账户和税费说明。组合层还需测试利率冲击与其他风险资产的共同回撤。没有这些资料，历史 0.59\% 年度 NAV 回报和 0.78\% 季度回报都不能外推。
\end{riskbox}

\section{对 3\% 现金目标的含义}

2025 年度现金分配只能证明该年度存在一次已披露分配，不能证明持有人未来可获得相同金额，更不能直接计入一个 3\% 的可复投净现金目标。只有实际结算、扣除税费和成本且法律性质已确认的现金，才可进入后续的净现金台账；其余市值或 NAV 变化只属于总回报的一部分。
